% Volume I: The Epistemology of Suspicion
% Part I. The Power of Paranoia
% Chapter 2: Why Appearances Underdetermine Causes
% Spine: proof-spines/ch002-spine.md

\chapter{Why Appearances Underdetermine Causes}
\label{ch:ch002}


Appearances underdetermine causes because what is observed is a projection of a richer latent process:
\[
Y_t=\mathcal O(X_t,\eta_t),
\]
where \(X_t\) is the causally relevant state, \(\mathcal O\) is an observation channel, and \(\eta_t\) collects noise, occlusion, selection, framing, compression, and measurement error. What arrives as evidence is therefore already filtered by a medium and a circumstance of capture.

The inverse problem is generally non-identifiable:
\[
\mathcal O(X_1,\eta_1)=\mathcal O(X_2,\eta_2)
\quad\text{while}\quad
X_1\neq X_2.
\]
An authentic video, statement, or measurement can therefore be real without being causally decisive. Truncated recordings, ambiguous agency, hidden antecedents, misleading timestamps, and observational equivalence all show how distinct histories can survive the same surface appearance.

This chapter accordingly distinguishes falsified evidence from incomplete evidence, decontextualized evidence, and causally insufficient evidence. Public argument often collapses those four categories into the blunt opposition between the ``real'' and the ``fake,'' but the epistemic problem is broader: a datum may be genuine and still fail to determine what happened.
